\begin{enumerate}
    \item \textbf{Difetto geometrico}

        Fissata una retta, un origine, un verso e una lunghezza unitaria, allora ad ogni $r \in \mathbb{Q}$
        corrisponde un punto P sulla retta. Non vale però l'inversa: ad ogni punto P della retta non corrisponde
        necessariamente un numero razionale.

        \begin{proposition}
            \label{pro:rdq-due-irr}
            $\nexists x\in\mathbb{Q} \mid x^2=2$
        \end{proposition}

        \begin{demonstration}
            \textit{di} \proref{pro:rdq-due-irr}

            \[\Th \nexists x\in\mathbb{Q} \mid x^2=2\]

            \begin{list_2}
                \item Supponiamo per assurdo che
                    \[ \Hp \exists x\in\mathbb{Q} \mid x^2=2\]
                \item Allora \label{pas:gcd-div-uno}
                    \[
                        \exists m \in \mathbb{Z} , \exists n\in\mathbb{N} \;\; | \;\; x = \frac{m}{n} \;\;
                        \land \;\; GCD(\left|m\right|, \left|n\right|)\footnote{\textbf{GCD:} Greatest Common Divisor, ovvero il massimo comune
                        divisore tra due o più numeri interi. In italiano viene più convenzionalmente indicato
                        con la sigla MCD.} = 1
                    \]
                \item Abbiamo dunque che
                    \[
                        x = \frac{m}{n}
                    \]\vspace{0.05cm}\[
                        x^2 = \left(\frac{m}{n}\right)^2 = \frac{m^2}{n^2} = 2
                    \]\vspace{0.05cm}\[
                        m^2 = 2n^2
                    \]\vspace{0.05cm}\[
                        2 \mid m^2
                    \]\vspace{0.05cm}\[
                        2 \mid \left|m\right|
                    \]\vspace{0.05cm}\[
                        \exists h \in \mathbb{N}  \mid  \left|m\right| = 2h
                    \]\vspace{0.05cm}\[
                        m^2 = 2n^2 = 4h^2
                    \]\vspace{0.05cm}\[
                        n^2 = 2h^2
                    \]\vspace{0.05cm}\[
                        2 \mid n^2
                    \]\vspace{0.05cm}\[
                        2 \mid \left|n\right|
                    \]\vspace{0.05cm}\[
                        GCD\left(\left|m\right|,\left|n\right|\right) \geq 2
                    \]
                \item Quest'ultima disequazione contraddice il passaggio \pasref{pas:gcd-div-uno}, la negazione della tesi inziale
                    non può quindi che essere falsa. Dunque la tesi è necessariamente vera.
            \end{list_2}

            \qed
        \end{demonstration}

        \item \textbf{Difetto algebrico}
        
        Non tutti i polinomi a coefficenti razionali hanno radici razionali

        \begin{example} \leavevmode
            \begin{list_3}
            

                \item Prendiamo
                \[
                    P(x) = x^2 - 2
                \]

                \item Le radici di $P(x)$ devono soddisfare la seguente equazione
                \[ P(x) = 0 \implies x^2 - 2 = 0 \implies x^2 = 2 \]
                \item Ma per \proref{pro:rdq-due-irr}

                \[ \nexists x\in\mathbb{Q}\mid x^2=2 \]
            \end{list_3}
        \end{example}

    \item \textbf{Difetto analitico}
    
    Il difetto analitico di $\mathbb{Q}$ è enunciato dalla proposizione del seguente teorema.

    \begin{theorem}
        \label{teo:non-com-raz}
        \textbf{Teorema della non completezza di $\mathbb{Q}$}

        L'insieme dei numeri razionali $\mathbb{Q}$ non è completo perché non gode della proprietà dell'estremo superiore
        (e, analogamente, non gode nemmeno della proprietà dell'estremo inferiore).
    \end{theorem}
    
    \begin{corollary}
        \textit{di} \teoref{teo:non-com-raz}
        \[
            \exists \mathbb{A} \subseteq \mathbb{Q} \mid
            \begin{cases}
                \mathbb{A}\neq\emptyset \\
                \mathbb{A} \text{ superiormente limitato}\\
                \nexists sup\:\mathbb{A} \notin \mathbb{Q} 
            \end{cases}
        \]
    \end{corollary}

    \begin{demonstration}
        \textit{di} \teoref{teo:non-com-raz}
        \[ \Th  \teoref{teo:non-com-raz} \Rightarrow\]
        \[ \Rightarrow \Th 
            \exists \mathbb{A} \subseteq \mathbb{Q} \mid
            \begin{cases}
                \mathbb{A}\neq\emptyset \\
                \mathbb{A} \text{ superiormente limitato}\\
                \nexists sup\:\mathbb{A} \notin \mathbb{Q} 
            \end{cases}    
        \]
        \begin{enumerate}[label=(\roman*)]
            \item Per affermare la tesi è dunque sufficente dimostrare che il
            sottoinsieme di $\mathbb{Q}$, definito come
            $\mathbb{A} := \left\{ x \in \mathbb{Q}^+ \mid x^2 \leq 2 \right\}$ soddisfa le seguenti condizioni.
            \[
                \begin{cases}
                    \mathbb{A}\neq\emptyset \\
                    \mathbb{A} \text{ superiormente limitato}\\
                    \nexists sup\:\mathbb{A} \notin \mathbb{Q} 
                \end{cases}   
            \]
            \item Appurato che $1\in\mathbb{A}$ poiché $1\in\mathbb{Q}^+ \;\land\; 1^2 = 1 \leq 2$, allora
            \[
                \mathbb{A}\neq\emptyset.
            \]
            \item Dimostrato che $2$ è un maggiorante per $\mathbb{A}$,
            allora $\mathbb{A}$ è superiormente limitato.

            \begin{enumerate}[label=\arabic*.]
                \item Dimostriamo che $2$ è un maggiorante per $\mathbb{A}$, cioè
                \[ \Th \forall x\in\mathbb{A}, 2 \geq x. \]
                \item Supponiamo per assurdo che
                \[
                    \exists a \in\mathbb{A}, 2 < a.
                \]
                \item Allora
                \[ a > 2 \] \[ a^2 > 2^2 \] \[ a^2 > 4 \] \[ a \notin \mathbb{A}.\]
                \item Viene contradetta la negazione della tesi per cui $a \in \mathbb{A}$.
                \item La negazione della tesi è necessariamente falsa.
                \item La tesi è necessariamente vera.
            \end{enumerate}
            \qed

            \item Mostrato che l'insieme $\mathbb{A}$ così definto è non vuoto e superiromente limitato,
            l'unica condizione rimanente è che definito $\Lambda := \sup\mathbb{A}$, allora
            \[ \nexists \Lambda := \sup\mathbb{A} \in \mathbb{Q}.\]

            \item Supponendo per assurdo che $\Lambda := \sup\mathbb{A} \in\mathbb{Q}$, basta dunque mostrare $\Lambda^2 = 2$,
            poiché contradirebbe \proref{pro:rdq-due-irr}, secondo cui $\nexists x\in\mathbb{Q} \mid x^2 = 2$.

            \item Per fare ciò mostriamo separatamente che $\Lambda^2 \geq 2$ e che $\Lambda^2 \leq 2$; da cui
            poi dedurremo che $\Lambda^2 = 2$.\

            \item \label{pas:lam-mag-zer} A tale scopo mostriamo prima che $\Lambda \geq 1$
            
            \begin{enumerate}[label=\arabic*.]
            
                \item Supponiamo per assurdo che $\Lambda < 1$.
                \item Essendo che
                \[
                    \begin{cases}
                        \frac{4}{3}\in\mathbb{Q}^+ \\
                        \left(\frac{4}{3}\right)^2 = \frac{16}{9} < \frac{18}{9} = 2\\
                        \frac{4}{3} > \frac{3}{3} = 1 > \Lambda
                    \end{cases}\Rightarrow
                    \begin{cases}
                        \frac{4}{3}\in\mathbb{Q}^+ \\
                        \left(\frac{4}{3}\right)^2 < 2\\
                        \frac{4}{3} > \Lambda
                    \end{cases}\Rightarrow
                    \begin{cases}
                        \frac{4}{3}\in\mathbb{A}\\
                        \Lambda < \frac{4}{3}
                    \end{cases} \Rightarrow
                \]
                \[\Rightarrow
                        \exists x = \frac{4}{3}\in\mathbb{A}, \Lambda < x. \]\[
                \]

                \item $\Lambda$ non è un maggiorante per $\mathbb{A}$, dunque viene violata la definizione precedente
                    secondo cui $\Lambda := \sup\mathbb{A}$, che richiederebbe che $\Lambda$ sia un maggiorante per $\mathbb{A}$.

                \item La negazione della tesi porta a una contraddizione.
                \item La tesi deve essere necessariamente vera.

            \end{enumerate}
            \qed

            \item \label{pas:due-low-bou} Dimostriamo che con $\Lambda\in\mathbb{Q}$, $\Lambda^2 \geq 2$

            \begin{enumerate}[label=\arabic*.]
                \item Abbiamo quindi
                \[ \Hp \Lambda\in\mathbb{Q} \quad \Th \Lambda^2 \geq 2.\]

                \item Supponiamo per assurdo che
                \[ \Lambda^2 < 2.\]

                \item \label{pas:con-eps} Prendendo quindi arbitrariamente $\epsilon\in\mathbb{Q} \mid 0 < \epsilon < 1$, allora
                \[ \Lambda + \epsilon > \Lambda .\]

                \item Sapendo che 
                \[ \epsilon > 0 \land  \underbrace{\Lambda > 0}_{\pasref{pas:lam-mag-zer}} \implies \Lambda+\epsilon\geq0.\]
                \item \label{pas:lam-rel-eps}Allora
                \[ (\Lambda + \epsilon)^2 > \Lambda^2 \]
                \[ \Lambda^2 + 2\Lambda\epsilon + \epsilon^2 > \Lambda^2 \]

                \item \label{pas:lam-con-eps} Da $\epsilon < 1 \land  \epsilon > 0$, abbiamo che
                
                \[ \epsilon < 1, \]
                \[ \epsilon ^2 < \epsilon, \]
                \[(\Lambda^2 + 2\Lambda\epsilon)+\epsilon^2 < (\Lambda^2 + 2\Lambda\epsilon) + \epsilon, \]
                \[\Lambda^2 + 2\Lambda\epsilon+\epsilon^2 < \Lambda^2 + 2\Lambda\epsilon + \epsilon, \]
                \[(\Lambda+\epsilon)^2 < \Lambda^2 + \epsilon \left(2\Lambda + 1\right). \]

                \item Prendiamo
                \[ \epsilon = \frac{2-\Lambda^2}{2\Lambda+1} \]

                e verifichiamo che soddisfi le condizioni di cui al punto \pasref{pas:con-eps} .

                \begin{enumerate}[label=\arabic{enumiii}.\arabic*.]
                    \item
                    
                    \[ \Th \epsilon\in\mathbb{Q}\]

                    Essendo $\mathbb{Q}$ chiuso rispetto a somma e prodotto (\pref{p:raz-chi-som},
                    \pref{p:raz-chi-pro}) ed essendo inolte sempre garantita l'esistenza dell'elemento inverso
                    in $\mathbb{Q}' = \mathbb{Q} \setminus \{0\}$ (\pref{p:inv-q}), allora

                    \[ 
                        \begin{cases}
                            \Lambda \in \mathbb{Q} \\
                            \underbrace{\Lambda \geq 0}_{\pasref{pas:lam-mag-zer}} \\
                        \end{cases}\Rightarrow
                        \begin{cases}
                            \Lambda\cdot\Lambda \in \mathbb{Q} \\
                            2\cdot\Lambda \in \mathbb{Q} \\
                            2\cdot\Lambda  \geq 0 \\
                        \end{cases}\Rightarrow
                        \begin{cases}
                            \Lambda^2 \in \mathbb{Q} \\
                            2\Lambda \in \mathbb{Q} \\
                            2\Lambda  \geq 0 \\
                        \end{cases}\Rightarrow\]\[\Rightarrow
                        \begin{cases}
                            - \Lambda^2 \in \mathbb{Q} \\
                            2\Lambda + 1 \in \mathbb{Q} \\
                            2\Lambda + 1  \geq 1 \\
                        \end{cases}\Rightarrow
                        \begin{cases}
                            - \Lambda^2 + 2 \in \mathbb{Q} \\
                            2\Lambda + 1 \in \mathbb{Q}' \\
                        \end{cases}\Rightarrow
                        \begin{cases}
                            2 - \Lambda^2 \in \mathbb{Q} \\
                            \frac{1}{2\Lambda + 1} \in \mathbb{Q}' \\
                        \end{cases}\Rightarrow\]\[\Rightarrow
                        \begin{cases}
                            2 - \Lambda^2 \in \mathbb{Q} \\
                            \frac{1}{2\Lambda + 1} \in \mathbb{Q} \\
                        \end{cases}\Rightarrow
                            \left(2 - \Lambda^2\right) \frac{1}{2\Lambda + 1} \in \mathbb{Q} \Rightarrow
                            \frac{2 - \Lambda^2}{2\Lambda + 1} \in \mathbb{Q} \Rightarrow\]\[\Rightarrow
                            \epsilon\in\mathbb{Q}
                    \]
                    \qed

                    \item
                    \[ \Th \epsilon > 0 \]
                    
                    \[
                        \begin{cases}
                            \Lambda^2 < 2 \\
                            \underbrace{\Lambda \geq 0}_{\pasref{pas:lam-mag-zer}}
                        \end{cases} \Rightarrow
                        \begin{cases}
                            2 - \Lambda^2 > 0 \\
                            2\cdot\Lambda \geq 0
                        \end{cases} \Rightarrow
                        \begin{cases}
                            2 - \Lambda^2 > 0\\
                            2\cdot\Lambda + 1 \geq 1
                        \end{cases}\Rightarrow\]\[\Rightarrow
                        \begin{cases}
                            2 - \Lambda^2 > 0\\
                            2\Lambda + 1 > 0
                        \end{cases}\Rightarrow
                        \frac{ 2 - \Lambda^2}{2\Lambda + 1}>0\Rightarrow
                        \epsilon>0
                    \]
                    \qed

                    \item
                    \[ \Th \epsilon < 1 \]

                    \[
                        \epsilon<1 \Leftarrow
                        \epsilon - 1 < 0 \Leftarrow
                        \frac{2 - \Lambda^2}{2\Lambda + 1} - 1 < 0 \Leftarrow\]\[\Leftarrow
                        \frac{2 - \Lambda^2}{2\Lambda + 1} - \frac{2\Lambda + 1}{2\Lambda + 1} < 0 \Leftarrow
                        \frac{2 - \Lambda^2 - \left(2\Lambda + 1\right)}{2\Lambda + 1} < 0 \Leftarrow\]\[\Leftarrow
                        \frac{2 - \Lambda^2 - 2\Lambda - 1}{2\Lambda + 1} < 0 \Leftarrow
                        \frac{- \Lambda^2 - 2\Lambda + 1}{2\Lambda + 1} < 0
                    \]


                    Sapendo che $\Lambda \geq 0 \Rightarrow 2\Lambda \geq 0 \Rightarrow
                    2\Lambda + 1\geq 1 \Rightarrow 2\Lambda + 1> 0$,

                    \[
                        \frac{- \Lambda^2 - 2\Lambda + 1}{2\Lambda + 1} < 0 \Leftarrow
                        - \Lambda^2 - 2\Lambda + 1 < 0 \Leftarrow
                        \Lambda^2 + 2\Lambda - 1 > 0 \Leftarrow\]\[\Leftarrow
                        \Lambda < -1-\sqrt{2}\; \lor\; \Lambda > -1+\sqrt{2}\Leftarrow
                        \Lambda > -1+\sqrt{2}\Leftarrow\]\[\Leftarrow
                        \Lambda \geq 1 \Leftarrow
                        \pasref{pas:lam-mag-zer}
                    \]
                    \qed
                \end{enumerate}

                \item Verificato che $\frac{2-\Lambda^2}{2\Lambda+1}$ è un valido valore per $\epsilon$,
                allora riprendendo l'ultimo passaggio di \pasref{pas:lam-con-eps} .
                
                \[
                    \left(\Lambda+\epsilon\right)^2 < \Lambda^2 + \epsilon\left(2\Lambda+1\right),\]\[
                    \left(\Lambda+\epsilon\right)^2 < \Lambda^2 + \frac{2 - \Lambda^2}{2\Lambda+1}\left(2\Lambda+1\right),\]\[
                    \left(\Lambda+\epsilon\right)^2 < \Lambda^2 + \left(2 - \Lambda^2\right),\]\[
                    \left(\Lambda+\epsilon\right)^2 < \Lambda^2 + 2 - \Lambda^2,\]\[
                    \left(\Lambda+\epsilon\right)^2 < 2. \]\[
                \]

                \item Abbiamo quindi che
                \[
                    \begin{cases}
                        \Lambda^2 < \left(\Lambda+\epsilon\right)^2 \text{ (vd. pass. \pasref{pas:lam-rel-eps})}\\ 
                        \left(\Lambda+\epsilon\right)^2 < 2
                    \end{cases} \]\[
                    \Lambda^2 < \left(\Lambda+\epsilon\right)^2 < 2 \]\[
                    \left(\Lambda+\epsilon\right)\in\mathbb{A}\]\[
                    \exists a=\left(\Lambda+\epsilon\right) \in\mathbb{A} \mid \Lambda < a
                \]

                \item $\Lambda$ non è un maggiorante per $\mathbb{A}$.
                \item $\Lambda$ non può essere l'estremo superiore di $\mathbb{A}$, ma $\Lambda:=\sup\mathbb{A}$.
                \item La negazione della tesi $\Lambda^2 \geq 2$ porta a una contraddizione.
                \item La tesi deve essere necessariamente vera.


            \end{enumerate}
            \qed

            \item \label{pas:due-upp-bou} Dimostriamo che con $\Lambda\in\mathbb{Q}$, $\Lambda^2 \leq 2$

            \begin{enumerate}[label=\arabic*.]
                \item Abbiamo quindi
                \[ \Hp \Lambda\in\mathbb{Q} \quad \Th \Lambda^2 \leq 2.\]

                \item \label{pass:ass-lam-mag} Supponiamo per assurdo che
                \[ \Lambda^2 > 2.\]

                \item \label{pas:con-eps/2}Prendendo quindi arbitrariamente $\epsilon\in\mathbb{Q} \mid 0 < \epsilon < \Lambda$, allora
                \[
                    0 < \Lambda - \epsilon < \Lambda,\]\[
                    \left(\Lambda - \epsilon\right)^2 < \Lambda^2, \]\[
                    \Lambda^2 -2\Lambda\epsilon + \epsilon^2 < \Lambda^2, \]\[
                    \Lambda^2 - 2\Lambda\epsilon  < \Lambda^2 - 2\Lambda\epsilon + \epsilon^2 < \Lambda^2, \]\[
                    \Lambda^2 - 2\Lambda\epsilon  < \left(\Lambda - \epsilon\right)^2  < \Lambda^2.
                \]

                \item Prendiamo
                \[ \epsilon = \frac{\Lambda^2-2}{2\Lambda} \]

                e verifichiamo che soddisfi le condizioni di cui al punto \pasref{pas:con-eps/2} .

                \begin{enumerate}[label=\arabic{enumiii}.\arabic*.]
                    \item
                    
                    \[ \Th \epsilon\in\mathbb{Q}\]

                    Essendo $\mathbb{Q}$ chiuso rispetto a somma e prodotto (\pref{p:raz-chi-som},
                    \pref{p:raz-chi-pro}) ed essendo inolte sempre garantita l'esistenza dell'elemento inverso
                    in $\mathbb{Q}' = \mathbb{Q} \setminus \{0\}$ (\pref{p:inv-q}), allora

                    \[ 
                        \begin{cases}
                            \Lambda \in \mathbb{Q} \\
                            \underbrace{\Lambda > 0}_{\pasref{pas:lam-mag-zer}} \\
                        \end{cases}\Rightarrow
                        \begin{cases}
                            \Lambda\cdot\Lambda \in \mathbb{Q} \\
                            2\cdot\Lambda \in \mathbb{Q} \\
                            2\cdot\Lambda  > 0 \\
                        \end{cases}\Rightarrow
                        \begin{cases}
                            \Lambda^2 \in \mathbb{Q} \\
                            2\Lambda \in \mathbb{Q} \\
                            2\Lambda  > 0 \\
                        \end{cases}\Rightarrow\]\[\Rightarrow
                        \begin{cases}
                            \Lambda^2 - 2 \in \mathbb{Q} \\
                            2\Lambda \in \mathbb{Q}'
                        \end{cases}\Rightarrow
                        \begin{cases}
                            \Lambda^2 - 2 \in \mathbb{Q} \\
                            \frac{1}{2\Lambda} \in \mathbb{Q}'
                        \end{cases}\Rightarrow
                        \begin{cases}
                            \Lambda^2 - 2 \in \mathbb{Q} \\
                            \frac{1}{2\Lambda} \in \mathbb{Q}
                        \end{cases}\Rightarrow\]\[\Rightarrow
                        \left(\Lambda^2 - 2 \right) \frac{1}{2\Lambda} \in \mathbb{Q} \Rightarrow
                        \frac{2 - \Lambda^2}{2\Lambda} \in \mathbb{Q} \Rightarrow
                        \epsilon\in\mathbb{Q}
                    \]
                    \qed

                    \item
                    \[ \Th \epsilon > 0 \]
                    
                    \[
                        \begin{cases}
                            \Lambda^2 > 2 \\
                            \underbrace{\Lambda > 0}_{\pasref{pas:lam-mag-zer}}
                        \end{cases} \Rightarrow
                        \begin{cases}
                            \Lambda^2 - 2> 0 \\
                            2\Lambda > 0
                        \end{cases} \Rightarrow
                        \frac{\Lambda^2 - 2}{2\Lambda}>0\Rightarrow
                        \epsilon>0
                    \]
                    \qed

                    \item
                    \[ \Th \epsilon < \Lambda \]

                    \[
                        \epsilon<\Lambda \Leftarrow
                        \epsilon - \Lambda < 0 \Leftarrow
                        \frac{\Lambda^2-2}{2\Lambda} - \Lambda < 0 \Leftarrow\]\[\Leftarrow
                        \frac{\Lambda^2-2}{2\Lambda} - \Lambda\frac{2\Lambda}{2\Lambda} < 0 \Leftarrow
                        \frac{\Lambda^2-2}{2\Lambda} - \frac{2\Lambda^2}{2\Lambda}< 0 \Leftarrow\]\[\Leftarrow
                        \frac{\Lambda^2 - 2 - 2\Lambda^2}{2\Lambda} < 0 \Leftarrow
                        \frac{- \Lambda^2 - 2}{2\Lambda} < 0
                    \]


                    Sapendo che $\Lambda > 0 \Rightarrow 2\Lambda > 0$,

                    \[
                        \frac{- \Lambda^2 - 2}{2\Lambda} < 0 \Leftarrow
                        - \Lambda^2 - 2 < 0 \Leftarrow
                        \Lambda^2 + 2 > 0 \Leftarrow\]\[\Leftarrow
                        \Lambda^2 > -2 \Leftarrow
                        \Lambda > 2\Leftarrow
                        (\pasref{pass:ass-lam-mag})
                    \]
                    \qed
                \end{enumerate}

                \item Verificato che $\frac{\Lambda^2-2}{2\Lambda}$ è un valido valore per $\epsilon$,
                allora riprendendo l'ultimo passaggio di \pasref{pas:con-eps/2} .
                
                \[
                    \Lambda^2 - 2\Lambda\epsilon  < \left(\Lambda - \epsilon\right)^2,\]\[
                    \Lambda^2 - 2\Lambda\frac{\Lambda^2-2}{2\Lambda}  < \left(\Lambda - \epsilon\right)^2,\]\[
                    \Lambda^2 - \left(\Lambda^2-2\right)  < \left(\Lambda - \epsilon\right)^2,\]\[
                    \Lambda^2 - \Lambda^2 + 2  < \left(\Lambda - \epsilon\right)^2,\]\[
                    2  < \left(\Lambda - \epsilon\right)^2.
                \]

                \item Grazie a quest'ultima disequazione mostriamo che $(\Lambda-\epsilon)$
                    deve essere un maggiorante per $\mathbb{A}$:
                \[
                    \forall a\in\mathbb{A}, a^2 \leq 2, \]\[
                    \forall a\in\mathbb{A}, a^2 \leq 2 < \left(\Lambda-\epsilon\right)^2, \]\[
                    \forall a\in\mathbb{A}, a^2 < \left(\Lambda-\epsilon\right)^2.
                \]

                Sapendo che $\epsilon<\Lambda \Rightarrow  0 < \Lambda-\epsilon
                \Rightarrow \Lambda-\epsilon > 0$, allora
                \[
                    \forall a\in\mathbb{A}, a < \left(\Lambda-\epsilon\right), \]\[
                    \forall a\in\mathbb{A}, \left(\Lambda-\epsilon\right) > a. 
                \]

                \item $(\Lambda-\epsilon)$ è un maggiorante per $\mathbb{A}$.
                \item $\Lambda$ non può essere l'estremo superiore di $\mathbb{A}$, ma $\Lambda:=\sup\mathbb{A}$.
                \begin{enumerate}[label=\arabic{enumiii}.\arabic*.]
                    \item Se $\Lambda$ fosse l'estremo superiore di $\mathbb{A}$, sarebbe anche necessariamente un
                    maggiorante per $\mathbb{A}$. Ma essendo anche $(\Lambda-\epsilon)$ un maggiorante per $\mathbb{A}$,
                    allora $\Lambda$ non sarebbe il minore dei maggioranti, perché $(\Lambda-\epsilon) < \Lambda$.
                    \qed
                \end{enumerate}
                \item La negazione della tesi $\Lambda^2 \leq 2$ porta a una contraddizione.
                \item La tesi deve essere necessariamente vera.


            \end{enumerate}
            \qed

            \item Abbiamo dunque da \pasref{pas:due-low-bou} e \pasref{pas:due-upp-bou} che,
            definendo $\Lambda:=\sup\mathbb{A}$,
            \[
                \begin{cases}
                    \Lambda^2 \geq 2 \\
                    \Lambda^2 \leq 2
                \end{cases}\implies \Lambda^2 = 2 \implies \nexists \Lambda \in \mathbb{Q}.
            \]

            \item La contraddizione della tesi $\nexists \Lambda \in \mathbb{Q}$, secondo cui
                $\exists \Lambda \in \mathbb{Q}$ porta a una contraddizione.

            \item La tesi $\nexists \Lambda \in \mathbb{Q}$ deve essere necessariamente vera.
            
            \item Esiste quindi un sottoinsieme di $\mathbb{Q}$ non vuoto e superiormente limitato che però non
                ammette un estremo superiore in $\mathbb{Q}$.

        \end{enumerate}
        \qed
    \end{demonstration}

\end{enumerate}